% Proposed B.137. New ordinary ambient-overlap result; no central changes.
\subsection*{The ordinary intersection of the two actual unfilled edge groups}

\statement{B.137}{Proved}{The two actual unfilled edge groups have trivial ordinary intersection, by exact vertex and incident-direction comparisons}
\label{stmt:B137}
Use the actual unfilled B.134/B.136 amalgam
$V=R^0*_{K}p(R^0)$, with $p|_K=\eta$ and its vertex-exchanging
involution $\theta$. Put $E_1=C\langle K,k\rangle C$ and
\begin{equation}\label{eq:B137-edge-groups}
 \begin{gathered}
 H_i=\langle E_1,s'\rangle
       =\langle E_1,a^{-1}p(k)a\rangle,\\
 H_j=\theta(H_i)=\langle p(E_1),s\rangle
       =\langle p(E_1),\eta(a)^{-1}k\eta(a)\rangle,\\
 B_i=\langle f,a^{-1}p(k)a\rangle\leq p(R^0),\qquad
 B_j=\langle\eta(f),\eta(a)^{-1}k\eta(a)\rangle\leq R^0.
 \end{gathered}
\end{equation}
The two displayed equalities retain the actual B.134 words
$w_0=aa_1^{-1}$, $v_0=p(w_0)$, $s'=w_0^{-1}p(k)w_0$,
$s=v_0^{-1}kv_0$, and $a_1=CaC$.
There are embedded ordinary subamalgams
\begin{equation}\label{eq:B137-subamalgams}
 H_i=E_1*_{\langle f\rangle}B_i,\qquad
 H_j=B_j*_{\langle\eta(f)\rangle}p(E_1).
\end{equation}
Their standard vertex and edge intersections are
\begin{equation}\label{eq:B137-standard-intersections}
 \begin{gathered}
 H_i\cap R^0=E_1,\quad H_i\cap p(R^0)=B_i,\quad
 H_i\cap K=\langle f\rangle,\\
 H_j\cap R^0=B_j,\quad H_j\cap p(R^0)=p(E_1),\quad
 H_j\cap K=\langle\eta(f)\rangle.
 \end{gathered}
\end{equation}
The actual vertex overlaps are trivial:
\begin{equation}\label{eq:B137-overlaps}
 E_1\cap B_j=1,\qquad p(E_1)\cap B_i=1.
\end{equation}
In fact their whole ordinary intersection is trivial:
\begin{equation}\label{eq:B137-whole-intersection}
 H_i\cap H_j=1.
\end{equation}
This conclusion uses the additional incident-direction comparison
proved below, not just the two standard vertex stabilizers.
No completed intersection, free-product assertion for the generated
overgroup, or survival after the two local fillings is inferred.

\begin{proof}
\emph{B.137.1. Exact based subamalgams.}
The element $a_1$ belongs to $E_1$, and
\[
 s'=a_1\bigl(a^{-1}p(k)a\bigr)a_1^{-1}.
\]
This proves the first generated-subgroup equality in
\eqref{eq:B137-edge-groups}. Applying $\theta$ and its actual
restriction $\eta$ on $K$ proves the second. Both conjugators
are therefore accounted for.

B.136 proves that $\langle\ell,p(k)\rangle$ is free on its
specified two generators, where $\ell=afa^{-1}$, and that its
intersection with $K$ is $\langle\ell\rangle$.
Conjugation by $a^{-1}\in K$ gives
\[
 B_i=\langle f\rangle*\langle a^{-1}p(k)a\rangle,
 \qquad B_i\cap K=\langle f\rangle.
\]
The other left-vertex factor has precisely that edge
intersection by B.136. Ordinary reduced-word normal form
therefore proves the first embedded subamalgam in
\eqref{eq:B137-subamalgams}. Its intersections with the two
ambient vertices and the edge are the first row of
\eqref{eq:B137-standard-intersections}. Applying $\theta$
proves the second subamalgam and second row, including the
actual $\eta$ on both edge words.

\emph{B.137.2. Reduce to the actual regional amalgam.}
Let $L=\langle K,k\rangle$ and $H_{\mathrm{ret}}=\langle h,k\rangle$,
where $h=fk$. The complete B.117 regional graph model has the
two color classes $\Omega_0,\Omega_1$ and gives
\[
 R^0=A(\Omega)\rtimes H_{\mathrm{ret}},\qquad
 L=A(\Omega_0)\rtimes H_{\mathrm{ret}},\qquad
 E_1=A(\Omega_1)\rtimes H_{\mathrm{ret}}.
\]
The last equality is literal: conjugation by $C$ exchanges the
colors and preserves $H_{\mathrm{ret}}$.
The color-zero retraction fixes $L$ and sends $E_1$ onto
$H_{\mathrm{ret}}$. Hence
\begin{equation}\label{eq:B137-color-intersection}
 E_1\cap L=H_{\mathrm{ret}}.
\end{equation}
Both generators of $B_j$ lie in $L$. It suffices to prove
$B_j\cap H_{\mathrm{ret}}=1$ inside the proper marked amalgam
\begin{equation}\label{eq:B137-regional-tree}
 \begin{gathered}
 L=H_{\mathrm{ret}}*_{T}P_{\mathrm{leaf}},\qquad T=\langle k\rangle,\\
 H_{\mathrm{ret}}=\langle h,k\mid[h^2,k]\rangle,\quad f=hk^{-1},\\
 P_{\mathrm{leaf}}=\langle k\rangle\times F(a,d,e,V_4),
 \qquad V_4=f^{-1}Bf.
 \end{gathered}
\end{equation}
This uses the entire regional graph and its actual retraction,
not a finite color or height sample.

\emph{B.137.3. The exact marked fork.}
Set
\begin{equation}\label{eq:B137-fork-words}
 \begin{gathered}
 \zeta=\eta(a),\qquad \ell=afa^{-1},\qquad q=\eta(\ell),
 \qquad \mathscr U=\langle k,q\rangle,\\
 \gamma_0=afV_4f^{-1}a^{-1}f,\qquad D_0=de d^{-1},\qquad
 w=\gamma_0D_0f^{-1}d^{-1}.
 \end{gathered}
\end{equation}
Then $\zeta B_j\zeta^{-1}=\mathscr U$.
The full certified five-word reflection gives the literal
free-word identities
\begin{equation}\label{eq:B137-two-branches}
 q=wf^{-1}w^{-1},\qquad
 \zeta=\gamma_0V_4^{-1}f^{-1}a^{-1}.
\end{equation}
The three $P_{\mathrm{leaf}}/T$ cosets represented by
\begin{equation}\label{eq:B137-directions}
 1,\qquad D_0,\qquad V_4^{-1}
\end{equation}
are distinct: their free-leaf coordinates are $1$, $de d^{-1}$,
and $V_4^{-1}$. These calculations retain the actual branch
and neck markings.

\emph{B.137.4. All-power endpoint tests and a tree immersion.}
Let $X$ be the ordinary Bass--Serre tree of
\eqref{eq:B137-regional-tree}. The element $k$ fixes the
$H_{\mathrm{ret}}$-vertex $aH_{\mathrm{ret}}$, because $a$
commutes with $k$. The element $q$ fixes the vertex
$wH_{\mathrm{ret}}$ by \eqref{eq:B137-two-branches}.
The reduced path $I$ between these vertices is given by
\[
 afV_4f^{-1}a^{-1}fD_0f^{-1}d^{-1}.
\]
Every displayed free-leaf block is nontrivial outside $T$,
and every $f^{\pm1}$ block is outside $T$ by the $h$-exponent
character of $H_{\mathrm{ret}}$. Thus $I$ is a genuine
geodesic between distinct vertices.

The first edge at $aH_{\mathrm{ret}}$ is $afT$, not $aT$.
For each nonzero integer $n$, equality of its translate by
$k^n$ with itself would require
\[
 f^{-1}a^{-1}k^na f=f^{-1}k^nf\in T.
\]
This cannot occur: the homomorphism
$H_{\mathrm{ret}}\to C_2(C)*\mathbb Z(k)$ sending $h$ to $C$
and fixing $k$ sends the word to $kCk^nCk^{-1}$, a reduced
word outside $\langle k\rangle$. Thus all incident endpoint
cosets supplied by powers of $k$ have distinct image edges.
At the other endpoint the last edge is $wT$; its translate
by $q^n$ can equal itself only if $f^{-n}\in T$, which the
$h$-exponent character excludes for every nonzero $n$.
This is the corresponding all-power endpoint test for $q$.

Start with the abstract free product
$\mathbb Z(\mathsf k)*\mathbb Z(\mathsf q)$ and subdivide
each edge of its standard tree into as many edges as $I$.
Map its fundamental subdivided edge to $I$, its endpoint
vertices to $aH_{\mathrm{ret}}$ and $wH_{\mathrm{ret}}$,
and extend equivariantly through the homomorphism sending
$\mathsf k,\mathsf q$ to the actual $k,q$. This is
well-defined since the endpoint factors fix the assigned
vertices. The all-power tests show local injectivity at
all translates of the endpoints. At subdivision vertices,
local injectivity is exactly the absence of backtracking
in $I$.

A locally injective map of simplicial trees is injective:
the image of every finite reduced path is reduced, so two
distinct vertices cannot have the same image. The standard
free-product action is faithful; the homomorphism of groups
is consequently injective as well. Let $Y\subseteq X$ be
the resulting embedded $\mathscr U$-invariant subtree.
Each interior subdivision vertex has degree two in $Y$
and trivial $\mathscr U$-stabilizer. These facts follow
from the actual tree embedding, not from an abstract rank
comparison.

\emph{B.137.5. The off-core branch and its stabilizer.}
The vertex $\gamma_0P_{\mathrm{leaf}}$ is an interior
subdivision vertex of the displayed fundamental segment.
Its two neighbors within $Y$ along that segment are
$\gamma_0H_{\mathrm{ret}}$ and
$\gamma_0D_0H_{\mathrm{ret}}$.
By \eqref{eq:B137-two-branches}, the path toward
$\zeta H_{\mathrm{ret}}$ begins instead through
$\gamma_0V_4^{-1}H_{\mathrm{ret}}$ and continues along
the reduced tail $f^{-1}a^{-1}$. Its first edge is distinct
from the two core directions by \eqref{eq:B137-directions}.

The degree of $Y$ at this vertex is exactly two, so the
third edge is outside $Y$. A subtree of a tree is convex.
The reduced path from this third edge cannot re-enter $Y$
without creating a cycle. Therefore the unique nearest
point of $\zeta H_{\mathrm{ret}}$ in $Y$ is
$\gamma_0P_{\mathrm{leaf}}$.

If $g\in\mathscr U$ fixes $\zeta H_{\mathrm{ret}}$, it
preserves $Y$ and fixes that unique nearest point. Its
stabilizer there is trivial by B.137.4, so $g=1$.
Equivalently
\[
 \mathscr U\cap\zeta H_{\mathrm{ret}}\zeta^{-1}=1,
 \qquad B_j\cap H_{\mathrm{ret}}=1.
\]
Combining with \eqref{eq:B137-color-intersection} proves
$E_1\cap B_j=1$. Applying $\theta$ gives
$p(E_1)\cap B_i=1$, proving \eqref{eq:B137-overlaps}.

\emph{B.137.6. The retraction controls every shared direction.}
The actual retraction $r:L\to K$ fixes $K$, kills $k$, and satisfies
\[
 r(H_{\mathrm{ret}})=\langle f\rangle\subseteq H_{\mathrm{ret}},
 \qquad r(B_j)=\langle\eta(f)\rangle\subseteq B_j.
\]
Take $b\in B_j\cap E_1K$ and write $b=e k_0$ with $e\in E_1$
and $k_0\in K$. Since $b,k_0\in L$, the element $e$ lies in
$E_1\cap L=H_{\mathrm{ret}}$; write it as $h_0$. Then
\[
 b\,r(b)^{-1}
   =h_0k_0\bigl(r(h_0)k_0\bigr)^{-1}
   =h_0r(h_0)^{-1}.
\]
The left side belongs to $B_j$, since $r(B_j)\subseteq B_j$.
The right side belongs to $H_{\mathrm{ret}}$.
By B.137.5 their intersection is trivial. Thus $b=r(b)\in K$,
and we have the exact identity
\begin{equation}\label{eq:B137-direction-intersection}
 B_j\cap E_1K=B_j\cap K=\langle\eta(f)\rangle.
\end{equation}
Consequently the incident edge-coset sets
$\{eK:e\in E_1\}$ and $\{bK:b\in B_j\}$ in $R^0/K$
have exactly the common member $K$. If $eK=bK$, then
$b\in E_1K$, and \eqref{eq:B137-direction-intersection}
forces $bK=K$. Applying the actual $\theta$ gives the same
statement on the right for $B_i$ and $p(E_1)$.
This controls directions, a strictly additional input beyond
the vertex intersections in \eqref{eq:B137-overlaps}.

\emph{B.137.7. The two invariant trees meet in exactly one edge.}
Let $\mathcal T$ be the ordinary tree of
$V=R^0*_{K}p(R^0)$, and let $e_0$ be its standard edge $K$.
Map the standard tree of the first subamalgam in
\eqref{eq:B137-subamalgams} into $\mathcal T$ by the actual
vertex maps. Its maps on incident edge cosets are injective
because $E_1\cap K=B_i\cap K=\langle f\rangle$.
Thus this tree map is locally injective, and hence globally
injective by the reduced-path argument of B.137.4.
Its image $\mathcal T_i$ is an embedded $H_i$-invariant
subtree containing $e_0$. The same argument, using
$B_j\cap K=p(E_1)\cap K=\langle\eta(f)\rangle$, produces
an embedded $H_j$-invariant subtree $\mathcal T_j$.

At the standard left vertex, their incident edges are
exactly the two coset sets in B.137.6. Their only shared
edge is $e_0$. At the right vertex the same is true by
$\theta$. The intersection of two nonempty subtrees of a
tree is a subtree, since the unique path between two common
vertices lies in both. If $\mathcal T_i\cap\mathcal T_j$
contained any edge or vertex beyond $e_0$, its path to
$e_0$ would depart from one endpoint along a shared edge.
Both possible departing directions have just been excluded.
Therefore
\[
 \mathcal T_i\cap\mathcal T_j=e_0,
\]
including its two endpoints.

Every element of $H_i\cap H_j$ preserves both invariant
trees, and hence this edge. The ambient amalgam action has
no edge inversions, so it fixes both endpoints. At the left
endpoint it lies in $E_1\cap B_j=1$. This proves the whole
ordinary intersection \eqref{eq:B137-whole-intersection}.
There is no unexamined cycle or shifted double coset in
this conclusion: the exact intersection of the actual
embedded trees supplies the global argument.

\emph{B.137.8. The remaining filling and compact scopes.}
The result does not identify the subgroup generated by
$H_i,H_j$ with their free product. Nor does it determine
intersections of their other conjugates, or normal-closure
interactions between their two local relator families.
Those issues remain relevant to survival of
$D_{\mathrm{loc}}$ in $V_*$ and to the $\kappa$-symmetry
of its actual kernel. The full two-sheet initial normal
subgroup $N_*$ and B.134's matching kernel remain unchanged.

All tree immersions, finite geodesics and nearest-point
arguments here are ordinary. No analogous completed
intersection or profinite branching result is asserted
without independent compact vertex and direction proofs.
\end{proof}

\begingroup\raggedright\footnotesize
The whole-input literal fork certificate is
\path{work_DE/next_marginal_source/check_ambient_vertex_fork28.py},
its genuine output, and
\path{work_DE/next_marginal_source/ambient_vertex_fork28.json}.
It reads the complete certified actual reflection dictionary;
the all-power tree argument is proved above. The complete
research proofs are
\path{work_DE/next_marginal_source/ambient_vertex_overlap28.md}
and
\path{work_DE/next_marginal_source/ambient_whole_overlap28.md}.
No finite subgroup or double-coset census is substituted for
the stated normal-form proof.
\par\endgroup
